%! Regression Lines and Residuals — Unit 2, Lesson 2.3 — AP Practice
% GRADUAL RELEASE: AP-format practice, assigned but NOT scored — the cover's score column reads
% NA for this component. The graded practice is the `homework` component that follows it.
% ONE PAGE: two multiple-choice items and one free-response set of three parts, all on one
% coffee-shop context (a fourth context — the notes use pizza deliveries and sleep, the
% homework used cars). MC 1 carries the sign-flip distractor (A); MC 2 is extrapolation.
% Free-response (c) is the spiral to Lesson 2.2 (direction, and association is not causation).
% Page lockstep with ap_practice_key rests on \writelines{n} in the blank matching n terse
% \ansline{} in the key. The next-lesson preview is NOT here: it closes the packet, in homework.
\documentclass[12pt]{article}
\usepackage{apstats-article}
\usepackage{apstats-boxes}

\begin{document}

\pageheader{Unit 2, Lesson 2.3}{AP Practice}
% NAMESTRIP: no name row here — it lives on the cover only.

\vspace{0.06in}

% Authored BYTE-IDENTICALLY in ap_practice/ and ap_practice_key/.
\boxguard[8]
\begin{remindbox}
\small
\textbf{Not required, not scored} --- AP reps. Your graded homework is the last section.
\end{remindbox}

\vspace{0.08in}

\boxguard[10]
\begin{scenariobox}[The Coffee Shop]{navy}
\small
A coffee shop recorded the daily high temperature ($^\circ$F) and the number of hot drinks
sold on 40 days with highs from 30$^\circ$F to 75$^\circ$F. The regression line is
\ $\widehat{\text{drinks}} = 210 - 2.4(\text{temperature})$.
\end{scenariobox}

\vspace{0.08in}

\boxguard[12]
\begin{headlinebox}{goldbg}
\small
\textbf{Section I --- Multiple Choice.} Circle the single best answer.
\end{headlinebox}

\begin{enumerate}[leftmargin=1.9em, itemsep=0.05in]

  \item On a 50$^\circ$F day the shop sold 100 hot drinks. What is the residual, and what
        does it mean?
        \begin{enumerate}[label=\textbf{(\Alph*)}, leftmargin=2.2em, itemsep=0pt]
          \item $-10$; the model predicted 10 more drinks than the shop sold.
          \item $10$; the shop sold 10 more hot drinks than the model predicted.
          \item $10$; the model predicted 10 more drinks than the shop sold.
          \item $90$; the model predicted 90 drinks, so the residual is 90.
          \item $-10$; the shop sold 10 more hot drinks than the model predicted.
        \end{enumerate}

  \item The manager plugs in a 100$^\circ$F day and gets $-30$ hot drinks. Which is the best
        explanation?
        \begin{enumerate}[label=\textbf{(\Alph*)}, leftmargin=2.2em, itemsep=0pt]
          \item The equation must be wrong, because sales cannot be negative.
          \item On a 100$^\circ$F day the shop will certainly sell no hot drinks.
          \item 100$^\circ$F is outside the 30--75$^\circ$F data, so this extrapolation is
                not trustworthy.
          \item The residual for a 100$^\circ$F day will be $-30$.
          \item The line predicts well for any temperature once it is rounded to zero.
        \end{enumerate}

\end{enumerate}

\vspace{0.02in}

\boxguard[12]
\begin{headlinebox}{goldbg}
\small
\textbf{Section II --- Free Response.} Show your work, and answer \emph{in context}.
\end{headlinebox}

\begin{enumerate}[leftmargin=1.9em, itemsep=0.1in, resume]

  \item \textbf{(a)} Predict the number of hot drinks sold on a 60$^\circ$F day. Show your
        work.\\[4pt]
        \writelines{1}

        \vspace{4pt}
        \textbf{(b)} On one 60$^\circ$F day the shop sold 58 hot drinks. Calculate the residual
        for that day and interpret it in context.\\[4pt]
        \writelines{2}

        \vspace{4pt}
        \textbf{(c)} Describe the direction of the association. Does warm weather \emph{cause}
        lower sales? Explain.\\[4pt]
        \writelines{2}

\end{enumerate}

\end{document}
